\clearpage\thispagestyle{empty}
\pdfbookmark[0]{Quellen zum Ceres-Bericht vom März 1805}{ceres180503-sources}
\begin{center}\textbf{Quellen zum Ceres-Bericht vom März 1805}\end{center}
\noindent Diese KI-Übertragung in heutiges Deutsch folgt dem Auszug aus Franz Xaver von Zachs Bericht in \href{https://gdz.sub.uni-goettingen.de/id/PPN236007467}{Gauß’ \textit{Werken}, Band~VI} (Göttingen, 1874), S.~261--262. Die entsprechenden Passagen in \textit{Ueber die Ceres}, Beitrag XXVI der \textit{Monatlichen Correspondenz}, Band~XI, März 1805, stehen auf \href{https://www.digitale-sammlungen.de/en/view/bsb10538603?page=295}{S.~283} und \href{https://www.digitale-sammlungen.de/en/view/bsb10538603?page=298}{S.~286}--\href{https://www.digitale-sammlungen.de/en/view/bsb10538603?page=299}{287}.
\par\smallskip\noindent Carl Friedrich Gauß berechnete die zehnten Elemente. Friedrich Wilhelm Bessel lieferte die für die Ephemeride benötigten Sonnenpositionen. Von Zach verweist auf Bessels \textit{Berechnung der Harriot’schen und Torporley’schen Beobachtungen des Cometen von 1607} im Novemberheft von 1804, \href{https://www.digitale-sammlungen.de/en/view/bsb10538602?page=439}{S.~425} und \href{https://www.digitale-sammlungen.de/en/view/bsb10538602?page=440}{folgende}.
\par\smallskip\noindent Übertragung und Satz wurden mit KI erstellt. Die Digitalisate der Zeitschrift stammen von der Bayerischen Staatsbibliothek.
\clearpage
